% Einheit 1 von 4 – Wurzelziehen und Lösbarkeit (bank/quadratische-gleichungen/e1.jsonl, zusammenbau v0.1)
%% TODO Zweigzeile Teil 1: was hier gelernt wird (Satz in Schülersprache aus dem Einheitstitel) – Einheit 1
\einheitenkopf[e1]{Einheit 1 von 4 · Wurzelziehen und Lösbarkeit}
\zweigzeile{ab Kl. 9, je nach Buch bis Kl. 10 (am Gymnasium ab Kl. 8) · P10}
\setcounter{aufgabe}{10}

%% TODO Titel: Ich-kann-Satz fehlt in der Bank (Platzhalter: Kettenname) – Nr. 11
%% TODO Anweisung: ein Satz über der ersten Teilaufgabe fehlt in der Bank – Nr. 11
\begin{aufgabe}{Quadratisch oder linear?}
\begin{teile}
\teil $5x - 2 = x^2$ – quadratisch oder linear?\\ \kreuz{quadratisch} \\ \kreuz{linear}
\end{teile}
\end{aufgabe}

%% TODO Titel: Ich-kann-Satz fehlt in der Bank (Platzhalter: Kettenname) – Nr. 12
%% TODO Anweisung: ein Satz über der ersten Teilaufgabe fehlt in der Bank – Nr. 12
\begin{aufgabe}{Wurzelziehen und Lösbarkeit}
%% TODO \rechenplatz{n}: Schrittzahl des Grundfalls fehlt in der Bank (nur bei fester Schreibform, 2.3 b)
\begin{teile}
\teil $(x - 3)^2 = 25$ – wie viele Lösungen?\\ \kreuz{zwei} \\ \kreuz{eine} \\ \kreuz{keine}
\end{teile}
\begin{gleichungsraster}[2]
\gl{\text{x^2 = 81}} &
\end{gleichungsraster}
\begin{teile}
\teil Ein quadratisches Beet hat den Flächeninhalt $64\,\text{m}^2$. Seitenlänge? \leerfeld[m]
\end{teile}
\begin{gleichungsraster}[2]
\gl{\text{x^2 = 13}} &
\end{gleichungsraster}
\begin{teile}
\teil $x^2 = -100$ – Lösungen? Begründe.
\end{teile}
\begin{gleichungsraster}[2]
\gl{\text{3x^2 = 75}} &
\end{gleichungsraster}
\begin{teile}
\teil Eine zylinderförmige Regentonne fasst $200$ Liter ($200\,\text{dm}^3$) und ist $8$ dm hoch. Für das Volumen gilt $V = \pi \cdot r^2 \cdot h$ ($r$ Radius, $h$ Höhe). Radius auf zwei Stellen? r = \leerfeld[dm]
\end{teile}
\begin{gleichungsraster}[2]
\gl{\text{2x^2 + 5 = 37}} &
\gl{\text{(x - 6)^2 = 16}} \\
\gl{\text{(x + 10)^2 = 25}} &
\gl{\text{(x - 9)^2 = 0}} \\
\end{gleichungsraster}
\begin{teile}
\teil Ist $x = -7$ Lösung von $(x + 2)^2 = 25$? \janein
\teil Parabel $y = (x - 1)^2 - 2$ und Gerade $y = 2$: Wie viele Lösungen hat $(x - 1)^2 - 2 = 2$?

\begin{ksys}[xmin=-5,xmax=5,ymin=-3,ymax=6,ablesen] \parabel{1}{1}{-2}{} \gerade{0}{2}{} \end{ksys}
Anzahl der Lösungen: \leerfeld
\teil Gib eine Gleichung der Form $x^2 = c$ ($c$ eine Zahl) an, die keine Lösung hat.
\teil Wahr oder falsch? Gib außerdem eine quadratische Gleichung ohne Lösung an. \hfill (P10 2021 OS)\\ A: $(x - 6)^2 = 0$ hat genau eine Lösung. \kreuz{wahr} \kreuz{falsch}\\ B: $7$ und $11$ sind Lösungen von $(x + 9)^2 = 4$. \kreuz{wahr} \kreuz{falsch} Gleichung ohne Lösung: \leerfeld
\end{teile}
\end{aufgabe}

%% TODO Titel: Ich-kann-Satz fehlt in der Bank (Platzhalter: Kettenname) – Nr. 13
%% TODO Anweisung: ein Satz über der ersten Teilaufgabe fehlt in der Bank – Nr. 13
\begin{aufgabe}{Wurzelziehen und Lösbarkeit – Fehler finden, Begründen, Darstellungswechsel, Anwendung}
\begin{teile}
\teil Wo ist der Fehler? \rechnung{x^2 &= 225 \\ x &= 15}
\teil Begründe: $x^2 = 10$ hat zwei Lösungen.
\teil Zeichne die Gerade $y = 4$ zur Normalparabel $y = x^2$ ein. Lösungen von $x^2 = 4$?

\begin{ksys}[xmin=-3,xmax=3,ymin=-1,ymax=6] \parabel{1}{0}{0}{} \end{ksys}
x1 ≈ \leerfeld , x2 ≈ \leerfeld
\teil Ein Stein fällt von einer $80$ m hohen Brücke. Für den Fallweg gilt $s = 5 \cdot t^2$ ($s$ in m, $t$ in s). Nach wie vielen Sekunden schlägt er auf? \leerfeld[s]
\end{teile}
\end{aufgabe}
